\begin{proof}
\textbf{Chebyshev’s Sum Inequality.}  
Let  

\[
a_{1}\le a_{2}\le\cdots\le a_{n},\qquad 
b_{1}\le b_{2}\le\cdots\le b_{n}
\]

be real numbers.  We shall prove  

\[
n\sum_{i=1}^{n}a_{i}b_{i}\;\ge\;
\Bigl(\sum_{i=1}^{n}a_{i}\Bigr)
\Bigl(\sum_{i=1}^{n}b_{i}\Bigr),
\tag{∗}
\]

and describe the equality case.

\begin{enumerate}
\item \textbf{Base case and degenerate sequences.}  
If \(n=1\) the inequality reads \(a_{1}b_{1}\ge a_{1}b_{1}\), true with equality.  

If all \(a_i\) are equal, say \(a_i=a\) for all \(i\), then  

\begin{align*}
n\sum_{i=1}^{n}a_i b_i
&=a\sum_{i=1}^{n}b_i
   =\Bigl(\sum_{i=1}^{n}a_i\Bigr)\Bigl(\sum_{i=1}^{n}b_i\Bigr),
\end{align*}
so (∗) holds as an equality.  The same argument works when the \(b_i\) are all equal.  
Henceforth we assume \(n\ge2\) and that \emph{not both} sequences are constant.

\item \textbf{A convenient notation.}  
Define  

\[
S:=n\sum_{i=1}^{n}a_{i}b_{i}
   -\Bigl(\sum_{i=1}^{n}a_{i}\Bigr)
    \Bigl(\sum_{i=1}^{n}b_{i}\Bigr).
\]

Our goal is to show \(S\ge0\).

\item \textbf{Transforming \(S\) into a symmetric double sum.}  

The product of the two simple sums can be written as  

\[
\Bigl(\sum_{i=1}^{n}a_i\Bigr)\Bigl(\sum_{j=1}^{n}b_j\Bigr)
   =\sum_{i=1}^{n}\sum_{j=1}^{n}a_i b_j .
\]

Likewise  

\[
n\sum_{i=1}^{n}a_i b_i
   =\sum_{i=1}^{n}\sum_{j=1}^{n}a_i b_i ,
\]

since for each fixed \(i\) the inner sum over \(j\) contributes exactly \(n\) copies
of the term \(a_i b_i\).  Hence  

\begin{align*}
S&=
\sum_{i=1}^{n}\sum_{j=1}^{n}a_i b_i
-\sum_{i=1}^{n}\sum_{j=1}^{n}a_i b_j  \\[2mm]
&=\frac12\sum_{i=1}^{n}\sum_{j=1}^{n}
   \bigl(a_i b_i-a_i b_j-a_j b_i+a_j b_j\bigr) \\[2mm]
&=\frac12\sum_{i=1}^{n}\sum_{j=1}^{n}
   (a_i-a_j)(b_i-b_j). \tag{1}
\end{align*}

Equation (1) is the desired double‑sum representation of \(S\).

\item \textbf{Non‑negativity of each summand.}  
Take any ordered pair \((i,j)\) with \(i<j\).  
Because the sequences are non‑decreasing,  

\[
a_i\le a_j,\qquad b_i\le b_j,
\]

so  

\[
a_i-a_j\le0,\qquad b_i-b_j\le0,
\]

and consequently  

\[
(a_i-a_j)(b_i-b_j)\ge0 .
\]

When \(i>j\) the same product appears (the double sum in (1) is symmetric);
therefore \((a_i-a_j)(b_i-b_j)\ge0\) for \emph{every} pair \((i,j)\).

From (1) we obtain  

\[
S=\frac12\sum_{i,j}(a_i-a_j)(b_i-b_j)\;\ge\;0 .
\]

Re‑substituting the definition of \(S\) yields precisely inequality (∗).

\item \textbf{Equality condition.}  
Equality \(S=0\) occurs iff each summand in (1) vanishes, i.e.  

\[
(a_i-a_j)(b_i-b_j)=0\qquad\text{for all }1\le i,j\le n .
\]

Thus for every unordered pair \(\{i,j\}\) we must have either \(a_i=a_j\) or
\(b_i=b_j\).  The only way this can happen for all pairs simultaneously is
that one of the two sequences is constant:

\begin{itemize}
\item either \(a_1=a_2=\dots =a_n\),
\item or \(b_1=b_2=\dots =b_n\).
\end{itemize}

The case \(n=1\) trivially satisfies both alternatives.  Hence equality in
Chebyshev’s inequality holds exactly when one of the sequences is constant.

\item \textbf{Concluding remark.}  
The same inequality follows from the rearrangement inequality, which states
that for two similarly ordered sequences the sum \(\sum a_i b_{\sigma(i)}\) is
maximal when the permutation \(\sigma\) preserves the ordering (the identity
permutation).  Taking \(\sigma\) to be the identity gives Chebyshev’s sum
inequality directly.

\end{enumerate}
\end{proof}